%!TEX root = ../main.tex

\chapter{Einleitung - Liam}

\section{Motivation und Ausgangssituation}
Stellen Sie sich vor, Sie navigieren während der Stoßzeit durch einen belebten Bahnhof. Sie verlassen sich auf Schilder, um den Ausgang zu finden, prüfen Abfahrtstafeln und achten auf andere Pendler. Für weltweit etwa 285 Millionen sehbehinderte Menschen stellt dieses Szenario eine tägliche Herausforderung voller Barrieren dar. Während der klassische Langstock Hindernisse am Boden erkennt, versagt er bei der Identifikation von Personen oder Objekten auf Kopfhöhe. Bestehende elektronische Hilfsmittel sind oft aufgrund extrem hoher Kosten (> 2.500 Euro) oder einer Abhängigkeit von permanenter Cloud-Konnektivität für die breite Masse unzugänglich.

\section{Problemstellung}
Das zentrale Problem liegt in der mangelnden digitalen Inklusion aufgrund finanzieller und technologischer Hürden. Mobile Assistenzsysteme müssen eine Echtzeit-Verarbeitung sensibler Bild- und Audiodaten ermöglichen, ohne die Ergonomie des Nutzers durch hohes Gewicht oder Wärmeentwicklung zu beeinträchtigen. Zudem erfordern Cloud-basierte Ansätze eine ständige Internetverbindung, was die Zuverlässigkeit in geschlossenen Räumen (z. B. U-Bahnhöfen) mindert und den Datenschutz gefährdet.

\section{Zielsetzung}
Das Ziel des Projekts \glqq Ascenta\grqq\ ist die Entwicklung eines Prototyps einer smarten Assistenzbrille, die visuelle Umweltinformationen in akustisches Feedback übersetzt. Das System soll kostengünstig (Materialkosten < 150 Euro) realisiert werden, eine diskrete Bauweise aufweisen und durch lokale Datenverarbeitung (Offline-Betrieb) maximale Datensicherheit für den Nutzer garantieren.

\section{Lösungsansatz und Kernidee}
Die Lösung basiert auf einer dezentralen \glqq Split-Brain\grqq-Architektur. Die Brille (Arduino Nicla Vision) fungiert als smarter Sensor-Satellit und führt eine ressourcensparende Objekterkennung (FOMO) zur Gesichtserkennung aus. Die rechenintensive Logik sowie die Spracherkennung werden auf das Smartphone des Nutzers ausgelagert. Die Übertragung erfolgt über ein latenzarmes WiFi-TCP-Protokoll. Durch diesen Ansatz wird die Hardware an der Brille entlastet, während die volle Prozessorleistung moderner Smartphones für die Sprachverarbeitung mittels Vosk genutzt wird.